%%%%%%%%%%%%%%%%%%%%%%%%%%%%%%%%%%%%%%%%%%%%%%%%%%%%%%%%%%%%%%%
%
% Welcome to Overleaf --- just edit your LaTeX on the left,
% and we'll compile it for you on the right. If you open the
% 'Share' menu, you can invite other users to edit at the same
% time. See www.overleaf.com/learn for more info. Enjoy!
%
%%%%%%%%%%%%%%%%%%%%%%%%%%%%%%%%%%%%%%%%%%%%%%%%%%%%%%%%%%%%%%%
\documentclass[17pt, a0paper, landscape,  margin=0mm, innermargin=10mm, blockverticalspace=3mm, colspace=3mm, subcolspace=3mm]{tikzposter}
\usepackage{graphicx}
\usepackage{authblk}
\usepackage{color}
% \usepackage[all]{xy}

\usetheme{board}

%% Elena's favorite green (thanks, Fernando!)
\definecolor{ForestGreen}{RGB}{34,139,34}
\definecolor{BlueViolet}{RGB}{138,43,226}
\definecolor{Coquelicot}{RGB}{255, 56, 0}
\definecolor{Teal}{RGB}{2,132,130}
%Uncomment this if you want to show work-in-progress comments
\newcommand{\comment}[1]{{\bf \tt  {#1}}}
% Uncomment this if you don't want to show comments
%\newcommand{\comment}[1]{}
% for color names see https://www.overleaf.com/learn/latex/Using_colors_in_LaTeX
\newcommand{\emcomment}[1]{\textcolor{ForestGreen}{\comment{Elena: {#1}}}}

\newcommand{\jscomment}[1]{\textcolor{magenta}{\comment{James: {#1}}}}


\title{Developing a framework for text-based game development for beginners in Clojure}
\author[1]{James Swearingen}
\author[1]{Elena Machkasova (adviser)}
\date{\today}
\affil[1]{Division of Science and Mathematics, University of Minnesota, Morris}
\usetheme{Envelope}
\usecolorstyle{Australia}

\begin{document}

\makeatletter
\def\maketitle{\AB@maketitle}
\makeatother

\maketitle

\block{Abstract}{
Clojure is a relatively new functional programming language that is used in industry and is supported by a community of developers. It is also a promising language for a beginner classroom due to its simple syntax, focus on small functions and on recursive algorithms, and data that’s immutable by default. However, there are gaps that make its adoption challenging. Debugging in Clojure can be unintuitive, because Clojure errors are exceptions in the underlying Java language. They frequently contain terminology intended for experienced developers, leaving out important context.

Past work at Morris included the modifications to the error messages to make them understandable by beginners, but in order to perform usability testing with the target audience the error messages need to be incorporated into a simple development environment. This project is a continuation of an ongoing research, focusing on developing such an environment.

% My work to be presented at MICS is focused on developing a framework for students to create simple text-based interactive games. The framework handles modifications to the game state while allowing students to write functions to describe the game state transformations without directly using mutability. Additionally, I am working on a set of simple functions that allow students to test their code without using sophisticated testing engines. This environment is inspired by a similar environment in the Racket programming language successfully used in introductory programming classes, although at this point we are not including graphical features used in Racket.

% My poster will present examples of uses of the environment. It will also discuss implementation challenges and their solutions. 
%\emcomment{Needs to be shortened and changed to match this year's work}
}
\begin{columns}
\column{0.3}
        \block{Overview of Functional Programming and Clojure}{
%\tkcomment{Tristan's section}
% \jscomment{This definitely needs some fixes; I don't expect anyone to know what a Lisp is, or the JVM.}
    \begin{itemize}
	% \item \emcomment{Includes reasons to teach FP to beginners}
	\item Clojure is a simple functional programming language [1].
	\item Functional programming encourages solving problems with small building blocks called functions.
	\item Functional languages have a rich history of being introduced early in programming education [2].
    \item A functional language is designed for programming using function composition rather than directly changing variables in computer memory. 
%Functional languages are immutable, meaning that when a variable is created, it cannot be changed. 
% this phrasing feels strange but I'm not sure what else to call them
    % \item In Clojure, functions are first-class citizens, which means functions can be declared, manipulated, and passed as arguments like most other data types.
    %promotes a more declarative programming style. Which is where you tell it what you want to happen rather than giving it a list of steps to follow.
    \item An important type of function is one that always produces the same output for the same input, and has no other effects.
%    \item Clojure is a functional language
%Lisp-based language 
%that encourages these principles and is a great tool for introducing students to functional programming.
    \item Clojure has a simple prefix notation syntax for function calls and special forms that's consistent across the language: \texttt{(+ 2 3), (filter odd? [1 2 3])}.
    \item It uses an interactive Read-Evaluate-Print-Loop (REPL) environment.
%: the user types some code, the system evaluates it and prints the answer.
    \end{itemize}
}

 \block{Clojure Error Messages}{
%\tkcomment{Tristan's section}
% \jscomment{Should we go back and take a look at the previous version of this section? I think we only changed it because of MICS.}
    \begin{itemize}
    \item Error messages are messages produced by a program when something goes wrong. A good error message should say what the problem is, and ideally make it clear how it should be fixed.
    \item Clojure code is compiled into Java. It is a more complex language whose data types and features are not known to beginner programmers.
    \item By the time Clojure code is executed, it is Java code, so Clojure errors appear as Java errors, along with a location in Clojure and an associated message that describes the error.
    \item Because the errors are Java errors, but the code and location are written in Clojure, this discrepancy can be confusing, in addition to the fact that the messages themselves are frequently unhelpful.


    % \item Error messages include Java error types and a Java-phrased cause, as well as a Clojure location.
    \end{itemize}

}

\block{Our Project, Babel}{

    \begin{itemize}
        \item The goal of the project is to use Clojure in an introductory programming class. 
%One focus of our research is development of a tool called Babel to modify Clojure error messages [4,5].
        \item To address the error messages issues, Babel provides a way to directly replace error messages in the REPL with simplified, beginner-friendly versions. [4,5]
        \item Errors on function calls give detailed information about expected arguments: types, position, etc.
        % \item Babel categorizes error data via dictionary lookup and specifications on Clojure function requirements. 
        % \item We developed and designed these specifications to capture as much information for the end user as possible, which is absent in default Clojure messages.
    \end{itemize}
        Correct call: \\
    \texttt{user=>(even? 2)} \\
    \texttt{true} \\
    Default Clojure Error Messages in REPL: \\
    \texttt{user=>(even? 2 3)} \\
    \texttt{Execution error (ArityException) at user/eval2046 (REPL:1).} \\ 
     \texttt{Wrong number of args (2) passed to: clojure.core/even?} \\
As in the previous error, the exact call to the failing function is not stated clearly. The only reference we get is the \texttt{clojure.core/even?}. \\

    Babel Error Messages in REPL: \\
    \texttt{user=>(even? 2 3)} \\
    \texttt{Wrong number of arguments in (even? 2 3): the function even? expects one argument but was given two arguments.} \\
    \texttt{In Clojure interactive session on line 1.} \\
The specific function call is stated in the error message, and the statement about "wrong number of args" is instead phrased in clear English. 

\vspace*{0.145in}
}



\column{0.4}
 \block{Testing Setup and Game Development Engine for Beginners}{
% \emcomment{Add some explanation of how it's related to error messages work}
    \begin{itemize}
        \item With the completion of the error message replacement system, I have shifted the focus to other teaching tools for beginner programmers.
        \item Most recently, I have completed a set of functions for testing that are more straightforward and lightweight than Clojure's built-in testing system, inspired by Beginner Student version of Racket used in How to Design Programs curriculum [3]. Examples of their use are listed below.
        \item My current project is a framework for students to create simple games. It simulates mutability with atoms, so that the students only need to write a simple set of functions to create the game. 
        \item The only things a student needs to write are pure functions. The framework handles the atoms behind the scenes, and only requires a map of game info and functions.
    \end{itemize}
    }
 \block{Examples of Interactions}{
    % \jscomment{maybe modify some of this since I'm going to be demonstrating on my laptop?}
    Check functions: \\
   All three of these return strings, not boolean values. The returned string tells you whether the test passed or failed, as well as what was tested. \\ \\
\texttt{check-equal} is the simplest. It tells you whether two values are equal.\\
\texttt{=> (check.fns/check-equal 1 1)} \\
\texttt{"Test (= 1 1) passed"} \\
\texttt{=> (check.fns/check-equal 1 2)} \\
\texttt{"Test (= 1 2) failed"} \\ \\
\texttt{check-range} checks whether a number is between two other numbers, inclusive, in the order of number, min, max. \\
\texttt{=> (check.fns/check-range 5 0 10)} \\
\texttt{"Test (<= 0 5 10) passed"} \\
\texttt{=> (check.fns/check-range 5 0 4)} \\
\texttt{"Test (<= 0 5 4) failed"} \\ \\
\texttt{check-precision} checks whether two numbers are within a certain distance of each other, in order of expected, actual, precision. \\
\texttt{=> (check.fns/check-precision 1.1 1 0.1)} \\
\texttt{"Test passed: 1 is within 0.1 of 1.1"} \\
\texttt{=> (check.fns/check-precision 2 1 0.3)} \\
\texttt{"Test failed: 1 is not within 0.3 of 2"} \\
Notice how with the check-equal and check-range functions, I've recreated the format of the internal function call. \\
This was done to show students what the function is doing internally, while still allowing simple testing. \\
    Game framework: \\
    There are three different components to the game: 
    \begin{itemize}
    \item the framework I am writing,
    \item the code that a student would write for their specific game, and
    \item the experience of a player.
    \end{itemize}
    To create a game, student programmers need to give the framework a few things: \\
    The initial state, a string telling players what commands are valid, and a set of functions to update the game state and draw it at the end of each turn. \\
    This framework allows for any plain-text game to be implemented without dealing with the complex nature of atoms in Clojure. \\

    This example is a recreation of tic-tac-toe. \\
    \texttt{=> (check.game/game check.tic-tac-toe/game-map)} \\
    This line calls the \texttt{game} function with the argument of \texttt{game-map}, which is a map that describes the game. This is all you need to create a game with this system. \\
\texttt{Possible commands: tl, tm, tr, ml, mm, mr, bl, bm, br}
\texttt{(tl for Top Left, etc.)} \\
\texttt{press Control+D to quit} \\
\texttt{tl} \\
\texttt{0 \_ \_} \\
\texttt{\_ \_ \_ } \\
\texttt{\_ 1 \_ } \\
In this example I played in the top left square, and the "enemy" (which chooses randomly) played in the bottom middle square. 

% \vspace*{0.07in}
}



\column{0.3}
\block{Results and Current Work}{
%\jscomment{Challenges + future work}

    \begin{itemize}
        \item With the completion of the error message system, I completed the checking functions, with the understanding that they would be used in an introductory class in the future.
        \item I also implemented the game framework, allowing students to create their own games.
        \item The current plan is for this game to be a project in the introductory class.
        \item There will be two options for this project: a one-player game or a two-player game.
        \item In a one-player game, the \texttt{win?} function will only need to check if one player has won.
        \item In a two-player game, the \texttt{win?} function needs to check both players.
        \item \texttt{enemy-turn} will also need to be defined, to simulate the enemy's turn.
        \item Unfortunately, it is possible for a student game to enter into an infinite recursion, and there isn't anything in place to handle this at the moment.
        % \item One potential solution is to put the game on a thread that times out.  
        \item I'm currently looking into a Clojure system called \textit{futures} to do this.
    \end{itemize}
}


\block{Future Work}{
    \begin{itemize}
        \item The next step for this project is to test with beginner programmers in a classroom setting.
        \item This will be part of a larger research project, where I plan to write a system to automatically log relevant information about Babel's error messages and how helpful they are.
        \item This would include data such as:
        \item Which errors students experience when working with their code?
        \item Which errors occur repeatedly with no sign of being resolved, and which ones are solved quickly?
        \item I am looking into a function called \texttt{tap>}, which might be useful for this.
    \end{itemize}
}

\block{Acknowledgments}{
 Thanks to Morris Academic Partnership (MAP) for supporting this project.\\
}

\block{References}{
    \begin{itemize}
        \item {[1]} \emph{clojure.org}
	\item {[2]} Matthias Felleisen, Robert Bruce Findler, Matthew Flatt, and Shriram Krishnamurthi \textit{The Structure and Interpretation of the Computer Science Curriculum},  Journal of Functional Programming, 2004.
    \item{[3]} Matthias Felleisen, Robert Bruce Findler, Matthew Flatt, Shriram Krishnamurthi \textit{How to Design Programs}, Second Edition, https://htdp.org/2023-3-6/Book/index.html
	\item{[4]} Myeongjae Song and Elena Machkasova \textit{Improving Clojure Error Messages for Programming Novices with clojure.spec} Midwest Instruction and Computing Symposium (MICS) 2017
	\item{[5]} Charlot Shaw \textit{Implementing Novice Friendly Error Messages in Clojure} Midwest Instruction and Computing Symposium (MICS) 2018
    \end{itemize}

    \vspace*{0.195in}
}

%\block{}{
%\begin{tikzfigure}
%    \includegraphics[width=0.55\linewidth]{figures/UMM logo.png}
%\end{tikzfigure}
%}

\end{columns}


\end{document}
